\appendix
\section{Implementation details}
\label{app:implementation}

This appendix holds the implementation details behind three steps of Section~\ref{sec:methods}. They
are needed to reproduce the analysis but not to follow its results.

\subsection{Compound identifier resolution}
\label{app:ids}

Matching the approval list to the screen was not done on names, because approval lists use
international nonproprietary names, often of salts, while screens use development codes and free bases.
In addition, CTRPv2's compound annotation \cite{ctrpv2} and DrEval's response file \cite{dreval} name
\NameMismatch{} of the \NDrugs{} compounds differently, \NameMismatchFda{} of them approved or clinical. Four keys were therefore
tried and the successful one recorded per compound: both name spaces, the PubChem compound identifier,
and PubChem's parent-compound relation \cite{pubchem}, which links a salt to the dissolved free base and
alone accounted for \NParentMatch{} matches, among them imatinib and doxorubicin. The two internal files were joined
on the screen's own compound identifier, which matched all \NDrugs{}. Cisplatin required one explicit
correction: it is recorded under a non-standard name with a structure identifier for elemental
platinum, and was resolved on the screen's recorded structure and vendor entry.

\subsection{Gene symbols and the scGPT vocabulary}
\label{app:hgnc}

scGPT matches genes to its vocabulary by symbol, and the expression matrix carries an older annotation of
the HUGO Gene Nomenclature Committee (HGNC) than the vocabulary \cite{hgnc}. Renamed genes would therefore
have been discarded although the vocabulary holds them under their current name: \NRescuedHvg{} genes of
the \NHvg-gene variant and \NRescuedAll{} of the full-transcriptome variant, the latter carrying
\PctRescuedAll\,\% of each cell's transcriptome, and because of the per-cell binning every remaining gene
would have shifted bin as well. The conversion step therefore recorded each gene's current HGNC symbol,
and the embedding step used it where the distributed symbol is not in the vocabulary, which raised the
\NHvg-gene variant from \NVocabDirect{} to \NVocab{} embedded genes.

\subsection{Confound covariates and cached fits}
\label{app:confounds}

The adjusted $R^2$ was used because the unadjusted one is biased upward by the number of covariates by an
amount that shrinks with the number of cells, so one bar would differ between a line of \CellsMin{} cells and one
of \CellsMax{}. On simulated predictions unrelated to the covariates the unadjusted value read about 19\,\%
and the adjusted about $-1\,\%$. Depth and mitochondrial fraction cannot be computed from the CPM matrix,
from which library size has been divided out and in which only four of 13 mitochondrial genes are
HVGs, so both were computed from the study's unnormalised UMI matrix over all genes: a total over HVGs
alone is depth multiplied by a biological share, and a confounder carrying the signal under test could
veto a true positive \cite{scp542,luecken2019}.
The per-fold standardisation and PCA fits of Section~\ref{sec:representations} were cached per run,
keyed by the reproducible fold assignment, in single precision like the all-cells decomposition, which
was kept for description only.
